% Part: lambda-calculus
% Chapter: syntax
% Section: substitution

\documentclass[../../../include/open-logic-section]{subfiles}

\begin{document}

\olfileid[pt]{lam}{syn}{sub}
\olsection{Substituição}

\begin{explain}
Variáveis livres são referências a variáveis do ambiente, portanto faz
sentido usar um valor específico no lugar de uma variável livre. Por exemplo,
podemos querer substituir $f$ em $\lambd[x][f x]$ por um termo específico,
como a função identidade $\lambd[y][y]$. Isso resulta em
$\lambd[x][(\lambd[y][y]) x]$. O processo de substituir variáveis livres por
termos lambda é chamado de substituição.
\end{explain}

\begin{defn}[Substituição] \ollabel{defn:substitution}
  A \emph{substituição} de um termo $N$ no lugar de uma variável $x$ em um termo
  $M$, $\Subst{M}{N}{x}$, é definida indutivamente por:
  \begin{enumerate}
    \item $\Subst{x}{N}{x} = N$. \ollabel{defn:substitution-1} 
    \item $\Subst{y}{N}{x}  = y$ se $x \neq y$. \ollabel{defn:substitution-2} 
    \item $\Subst{PQ}{N}{x}  = (\Subst{P}{N}{x}) (\Subst{Q}{N}{x})$.
      \ollabel{def:substitution-3}
    \item $\Subst{(\lambd[y][P])}{N}{x} = \lambd[y][\Subst{P}{N}{x}]$,
      se $x \neq y$ e $y \notin \FV{N}$; caso contrário, não está definida.
      \ollabel{defn:substitution-4}
  \end{enumerate}
\end{defn}

\begin{explain}
Em \olref{defn:substitution}\olref{defn:substitution-4}, exigimos que $x
\neq y$ porque não queremos substituir as ocorrências \emph{ligadas} da
variável~$x$ em~$M$ por~$N$. Por exemplo, se calcularmos a
substituição $\Subst{\lambd[x][x]}{y}{x}$, o resultado não deve ser
$\lambd[x][y]$, mas simplesmente $\lambd[x][x]$.

Ao substituir $N$ no lugar de~$x$ em $\lambd[y][P]$, também exigimos que
$y \notin \FV{N}$. Por exemplo, não podemos substituir $y$ no lugar de $x$ em
$\lambd[y][x]$, isto é, $\Subst{\lambd[y][x]}{y}{x}$, porque isso resultaria
em $\lambd[y][y]$, termo que representa a função que recebe um argumento e o
devolve diretamente. Entretanto, o termo $\lambd[y][x]$ representa uma
função que sempre devolve o termo~$x$ (ou aquilo a que $x$ se refere).
Portanto, o resultado que de fato queremos é uma função que recebe um
argumento, descarta-o e devolve a variável de ambiente~$y$. Para fazer isso
corretamente, primeiro teríamos de «renomear» a variável ligada~$y$.
\end{explain}

\begin{prob}
  Qual é o resultado das substituições a seguir?
  \begin{enumerate}
  \item $\Subst{\lambd[y][x(\lambd[w][vwx])]}{(uv)}{x}$
  \item $\Subst{\lambd[y][x(\lambd[x][x])]}{(\lambd[y][xy])}{x}$
  \item $\Subst{y(\lambd[v][xv])}{(\lambd[y][vy])}{x}$
  \end{enumerate}
\end{prob}

\begin{thm} \ollabel{thm:notinfv}
  Se $x \notin \FV{M}$, então $\FV{\Subst{M}{N}{x}} = \FV{M}$, se
  o lado esquerdo estiver definido.
\end{thm}

\begin{proof}
  Por indução sobre a formação de $M$.
  \begin{enumerate}
  \item $M$ é uma variável: exercício.
  \item $M$ é da forma $(PQ)$: exercício.
  \item $M$ é da forma $\lambd[y][P]$ e, como
    $\Subst{\lambd[y][P]}{N}{x}$ está definida, ela deve ser
    $\lambd[y][\Subst{P}{N}{x}]$. Então $\Subst{P}{N}{x}$ deve estar
    definida; além disso, $x \neq y$ e $x \notin \FV{Q}$. Logo:
    \begin{multline*}
      \FV{\Subst{\lambd[y][P]}{N}{x}} = \\
    \begin{aligned}
      & = \FV{\lambd[y][\Subst{P}{N}{x}]} &&
      \text{por \olref{defn:substitution-4}}\\
      & = \FV{\Subst{P}{N}{x}} \setminus \{y\} &&
      \text{por \olref[fv]{def:fv}\olref[fv]{def:fv2}}\\
      & = \FV{P} \setminus \{y\} && \text{pela hipótese de indução} \\
      & = \FV{\lambd[y][P]} && \text{por \olref[fv]{def:fv}\olref[fv]{def:fv2}}
    \end{aligned}
    \end{multline*}
  \end{enumerate}
\end{proof}

\begin{prob}
  Complete a demonstração de \olref[lam][syn][sub]{thm:notinfv}.
\end{prob}

\begin{thm} \ollabel{thm:infv}
  Se $x \in \FV{M})$, então $\FV{\Subst{M}{N}{x}} = (\FV{M} \setminus
  \{x\}) \cup \FV{N}$, desde que o lado esquerdo esteja definido.
\end{thm}

\begin{proof}
  Por indução sobre a formação de $M$.
  \begin{enumerate}
  \item $M$ é uma variável: exercício.
  \item $M$ é da forma $PQ$: como
    $\Subst{(PQ)}{N}{y}$ está definida, ela deve ser
    $(\Subst{P}{N}{x})(\Subst{Q}{N}{x})$, com ambas as substituições
    definidas. Além disso, como $x \in \FV{PQ}$, temos $x \in \FV{P}$ ou
    $x \in \FV{Q}$, ou ambos. O restante fica como exercício.
  \item $M$ é da forma $\lambd[y][P]$. Como
    $\Subst{\lambd[y][P]}{N}{x}$ está definida, ela deve ser
    $\lambd[y][\Subst{P}{N}{x}]$, com $\Subst{P}{N}{x}$
    definida, $x \neq y$ e $y \notin \FV{N}$; além disso, como $y \in
    \FV{\lambd[x][P]}$, também temos $y \in \FV{P}$. Agora:
    \begin{multline*}
      \FV{\Subst{(\lambd[y][P])}{N}{x}} = \\
      \begin{aligned}
      & = \FV{\lambd[y][\Subst{P}{N}{x}]} \\
      & = \FV{\Subst{P}{N}{x}} \setminus \{y\} \\
      & = ((\FV{P} \setminus \{y\}) \cup (\FV{N} \setminus \{x\})
       && \text{pela hipótese de indução} \\
      & = (\FV{P} \setminus \{x, y\}) \cup \FV{N}
       && x \notin \FV{N} \\
      & = (\FV{\lambd[y][P]} \setminus \{x\}) \cup \FV{N}
      \end{aligned}
      \end{multline*}
  \end{enumerate}
\end{proof}

\begin{prob}
  Complete a demonstração de \olref[lam][syn][sub]{thm:infv}.
\end{prob}

\begin{thm}\ollabel{thm:clr}
  $x \notin \FV{\Subst{M}{N}{x}}$, se o lado direito estiver definido e
  $x \notin \FV{N}$.
\end{thm}

\begin{proof}
  Exercício.
\end{proof}

\begin{prob}
  Demonstre \olref[lam][syn][sub]{thm:clr}.
\end{prob}


\begin{thm}\ollabel{thm:inv}
  Se $\Subst{M}{y}{x}$ está definida e $y \notin \FV{M}$, então
  $\Subst{\Subst{M}{y}{x}}{x}{y} = M$.
\end{thm}

\begin{proof}
  Por indução sobre a formação de $M$.
  \begin{enumerate}
  \item $M$ é uma variável $z$: exercício.
  \item $M$ é da forma $(PQ)$. Então:
    \begin{align*}
      \Subst{\Subst{(PQ)}{y}{x}}{x}{y}
      &=\Subst{((\Subst{P}{y}{x})(\Subst{Q}{y}{x}))}{x}{y} \\
      &= (\Subst{\Subst{P}{y}{x}}{x}{y})(\Subst{\Subst{Q}{y}{x}}{x}{y}) \\
      &= (PQ) \text{ pela hipótese de indução}
    \end{align*}
  \item $M$ é da forma $\lambd[z][N]$. Como
    $\Subst{\lambd[z][N]}{y}{x}$ está definida, sabemos que
    $z \neq y$. Portanto:
    \begin{align*}
      \Subst{\Subst{(\lambd[z][N])}{y}{x}}{x}{y}\\
      & = \Subst{(\lambd[z][\Subst{N}{y}{x}])}{x}{y} \\
      & = \lambd[z][\Subst{\Subst{N}{y}{x}}{x}{y}] \\
      &= \lambd[z][N] \text{ pela hipótese de indução} 
    \end{align*}
  \end{enumerate}
\end{proof}

\begin{prob}
  Complete a demonstração de \olref[lam][syn][sub]{thm:inv}.
\end{prob}

\end{document}
